\documentclass[11pt,a4paper]{article}

% --- Essential Packages ---
\usepackage[utf8]{inputenc}
\usepackage[T1]{fontenc}
\usepackage[margin=1in]{geometry}
\usepackage{amsmath,amssymb,amsfonts}
\usepackage{booktabs}
\usepackage{graphicx}
\usepackage{xcolor}
\usepackage{hyperref}
\usepackage{fancyhdr}
\usepackage{microtype}
\usepackage{titlesec}
\usepackage{float}
\usepackage{setspace}
\usepackage{enumitem}

% --- Hyperlink Setup ---
\hypersetup{
    colorlinks=true,
    linkcolor=blue!70!black,
    citecolor=blue!70!black,
    urlcolor=blue!70!black
}

% --- Header and Footer ---
\pagestyle{fancy}
\fancyhf{}
\rhead{\small \textcolor{gray}{Finesse $\times$ Citadel $\mid$ Portfolio Strategy}}
\lhead{\small \textcolor{gray}{Systematic Quantitative Research}}
\cfoot{\thepage}

% --- Section Styling ---
\titleformat{\section}{\large\bfseries\color{blue!60!black}}{}{0em}{}[\titlerule]
\titleformat{\subsection}{\normalsize\bfseries\color{black!85}}{}{0em}{}
\titleformat{\subsubsection}{\small\bfseries\color{black!70}}{}{0em}{}

\setstretch{1.15}

\begin{document}

% --- Title Block ---
\begin{center}
    {\LARGE \textbf{Systematic Multi-Cap Equity Strategy with Dual-Factor Momentum, Dynamic De-Correlation, and Risk-Parity Allocation}}\\[0.8em]
    {\large \textbf{Finesse $\times$ Citadel Portfolio Construction Challenge (Round 2)}}\\[0.5em]
    \rule{\linewidth}{0.8pt}\\[0.4em]
    \textbf{Initial Corpus:} Rs.~1,00,00,000 (1.00 Crore) \quad $\mid$ \quad \textbf{Universe:} Nifty 100, Midcap 100, Smallcap 100 ($\sim$300 Stocks) \quad $\mid$ \quad \textbf{Period:} 2021--2025
\end{center}

\vspace{0.5em}

% ==============================================================================
% 1. EXECUTIVE SUMMARY & PROBLEM FORMULATION
% ==============================================================================
\section{1. Executive Summary \& Strategy Philosophy}
This document specifies a fully systematic, rule-based quantitative equity investment strategy designed for the Indian equity universe ($\sim$300 constituents of the Nifty 100, Nifty Midcap 100, and Nifty Smallcap 100 indices). Managing an initial virtual capital of Rs.~1,00,00,000 (Rs.~1 Crore) with a strict constraint of holding at most 10 stocks at any point in time, the model maximizes Total Net Profit and Loss (PNL) while embedding rigorous multi-layered risk controls to prevent catastrophic drawdowns and sector concentration.

The strategy operates on a core market premise: asset prices exhibit intermediate-term trend persistence driven by gradual institutional capital reallocation, but unconstrained momentum portfolios frequently suffer severe drawdowns due to hidden cross-asset correlations and high-beta volatility. To solve this, our model combines:
\begin{enumerate}[leftmargin=1.8em, itemsep=0.1em]
    \item A long-term trend filter to eliminate assets in secular decline.
    \item A dual-factor composite ranking that balances intermediate trend strength with price trajectory stability.
    \item A greedy pairwise de-correlation filter to enforce true statistical diversification.
    \item An inverse-volatility (risk-parity) allocation engine that equalizes risk contributions across market-cap tiers while strictly accounting for 0.1\% per-trade transaction friction.
\end{enumerate}

% ==============================================================================
% 2. FACTOR DEFINITIONS, MATHEMATICAL FORMULATIONS & ECONOMIC LOGIC
% ==============================================================================
\section{2. Factor Framework: Definitions, Math, and Financial Logic}
The strategy evaluates each stock in the 300-asset universe across four distinct quantitative components. Below is the detailed breakdown of each factor, its exact formula, and the economic rationale behind its selection.

\subsection{2.1 Trend Regime Filter: 200-Day Simple Moving Average (SMA)}
\begin{itemize}[leftmargin=1.5em]
    \item \textbf{Mathematical Definition:} Let $P_{i, t}$ be the adjusted closing price of stock $i$ on date $t$. The 200-day Simple Moving Average is computed as:
    \begin{equation}
        \text{SMA}_{200}(P_{i, t}) = \frac{1}{200}\sum_{\tau=0}^{199} P_{i, t-\tau}
    \end{equation}
    The eligibility condition requires:
    \begin{equation}
        P_{i, t} > \text{SMA}_{200}(P_{i, t})
    \end{equation}
    \item \textbf{Financial Logic \& Intuition:} The 200-day SMA is a classical institutional indicator separating secular bull phases from secular bear regimes. Stocks trading below their 200-day average often suffer from deteriorating business fundamentals, earnings downgrades, or institutional distribution. Enforcing this filter prevents the strategy from buying ``falling knives'' or catching structural value traps.
\end{itemize}

\subsection{2.2 Factor 1: 12-Minus-1 Month Residual Price Momentum ($M_{i,t}$)}
\begin{itemize}[leftmargin=1.5em]
    \item \textbf{Mathematical Definition:} For each eligible stock $i$, momentum is calculated over a 250-trading-day ($\sim$1 year) lookback, explicitly skipping the most recent 21 trading days ($\sim$1 month):
    \begin{equation}
        M_{i, t} = \frac{P_{i, t-21}}{P_{i, t-250}} - 1
    \end{equation}
    \item \textbf{Financial Logic \& Intuition:} 
    \begin{enumerate}[label=(\alph*), itemsep=0.1em]
        \item \textit{The 12-Month Momentum Effect:} Financial markets exhibit persistent drift over 3- to 12-month horizons. This occurs because institutional investors accumulate large positions gradually over quarters, and market participants systematically underreact to sustained fundamental improvements (such as successive quarters of earnings growth).
        \item \textit{Why Skip the Last 21 Days (The ``Minus-1'' Month):} Short-term price returns (over 1 to 4 weeks) are heavily contaminated by liquidity shocks, bid-ask bounce, and mean-reversion. Stocks that surged aggressively over the last 10 days often experience immediate short-term pullbacks. By measuring the price movement from day $t-250$ to day $t-21$, we isolate true intermediate-term trend persistence while filtering out short-term noise.
    \end{enumerate}
\end{itemize}

\subsection{2.3 Factor 2: Realized Volatility \& Trajectory Stability ($V_{i,t}$)}
\begin{itemize}[leftmargin=1.5em]
    \item \textbf{Mathematical Definition:} Realized volatility ($\sigma_{i,t}$) is computed as the sample standard deviation of daily percentage returns $r_{i,\tau} = \frac{P_{i,\tau} - P_{i,\tau-1}}{P_{i,\tau-1}}$ over the trailing 250 trading days, annualized by $\sqrt{252}$:
    \begin{equation}
        \sigma_{i, t} = \sqrt{252} \cdot \sqrt{\frac{1}{T-1}\sum_{\tau=t-250}^{t} (r_{i,\tau} - \bar{r}_i)^2}
    \end{equation}
    The stability factor $V_{i,t}$ is defined as the inverse realized volatility:
    \begin{equation}
        V_{i, t} = \frac{1}{\sigma_{i, t} + 10^{-6}}
    \end{equation}
    \item \textbf{Financial Logic \& Intuition:} In quantitative asset pricing, the \textit{Low-Volatility Anomaly} demonstrates that assets with smoother, low-variance return profiles generate superior risk-adjusted compounded returns over long horizons. Wildly fluctuating stocks often experience steep drawdown cycles that destroy geometric compounding. Favoring higher stability ($V_{i,t}$) rewards steady, compounding price appreciation over speculative, lottery-like volatility spikes.
\end{itemize}

\subsection{2.4 Cross-Sectional Normalization \& Composite Alpha Formulation ($S_i$)}
Because raw momentum is bounded as a percentage return while inverse volatility is in inverse standard deviation units, both signals are standardized cross-sectionally across the eligible universe $\mathcal{U}_t$ on each rebalance date:
\begin{equation}
    Z_{\text{mom}, i} = \frac{M_{i,t} - \mu_M}{\sigma_M}, \qquad Z_{\text{vol}, i} = \frac{V_{i,t} - \mu_V}{\sigma_V}
\end{equation}
where $\mu_M, \sigma_M$ and $\mu_V, \sigma_V$ represent the cross-sectional mean and standard deviation of $M$ and $V$ across all eligible stocks on day $t$.

The final \textbf{Composite Alpha Score} $S_i$ is computed as a weighted linear combination:
\begin{equation}
    S_i = 0.60 \cdot Z_{\text{mom}, i} + 0.40 \cdot Z_{\text{vol}, i}
\end{equation}
\textit{Weighting Rationale:} Giving a 60\% weight to momentum ensures robust capture of strong upward price trends, while the 40\% stability weight acts as a quality damper, penalizing erratic and bubble-driven assets.

% ==============================================================================
% 3. PORTFOLIO SELECTION & RISK DE-CORRELATION
% ==============================================================================
\section{3. Stock Selection \& Dynamic De-Correlation Rules}
Simply picking the top 10 highest-scoring stocks often leads to severe unintended concentration (e.g., all 10 stocks belonging to a single hot theme like PSU energy or IT). To enforce true structural diversification, candidate stocks sorted descending by $S_i$ must pass two greedy screening filters:

\subsection{3.1 Sector Concentration Cap}
No more than \textbf{2 stocks} may belong to the same industry classification:
\begin{equation}
    \text{Count}\left(\{k \in \mathcal{P} \mid \text{Sector}(k) = \text{Sector}(i)\}\right) < 2
\end{equation}
This prevents industry-specific regulatory or macroeconomic shocks from impairing the aggregate portfolio.

\subsection{3.2 Rolling 60-Day Pairwise De-Correlation Screen}
For any candidate stock $i$, its rolling 60-day Pearson correlation $\rho_{ij}$ is evaluated against every security $j$ already accepted into the target portfolio $\mathcal{P}$:
\begin{equation}
    \rho_{ij} = \frac{\sum_{\tau=t-60}^{t} (r_{i,\tau} - \bar{r}_i)(r_{j,\tau} - \bar{r}_j)}{\sqrt{\sum_{\tau=t-60}^{t}(r_{i,\tau}-\bar{r}_i)^2} \sqrt{\sum_{\tau=t-60}^{t}(r_{j,\tau}-\bar{r}_j)^2}}
\end{equation}
The candidate stock is admitted into the portfolio if and only if:
\begin{equation}
    \max_{j \in \mathcal{P}} \rho_{ij} \le 0.65
\end{equation}
\textit{Financial Logic:} Even across different sectors, equities often exhibit high co-movement during market regimes. Rejecting candidate stocks with pairwise correlation $\rho > 0.65$ ensures that every asset added contributes independent variance, lowering overall portfolio risk. Selection proceeds down the ranked list until $K = 10$ assets are selected.

% ==============================================================================
% 4. CAPITAL ALLOCATION & WEIGHTING SCHEME
% ==============================================================================
\section{4. Capital Allocation \& Risk-Parity Weighting}
Once the $K \le 10$ constituents are selected, capital is distributed using an \textbf{Inverse-Volatility Allocation (Risk-Based Weighting)} rather than naive equal weighting:

\begin{equation}
    w_i = \frac{\frac{1}{\sigma_{i, t}}}{\sum_{j=1}^K \frac{1}{\sigma_{j, t}}} \quad \text{such that} \quad \sum_{i=1}^K w_i = 1.0
\end{equation}

\subsection{4.1 Intuitive Rationale Behind Risk-Based Allocation}
In a multi-cap portfolio containing both large-cap blue chips (volatility $\approx 15\%$) and small-cap growth names (volatility $\approx 40\%$), equal weighting (10\% capital per stock) results in an unbalanced portfolio where small caps generate over 75\% of daily portfolio risk. 

Under inverse-volatility weighting:
\begin{itemize}[leftmargin=1.5em, itemsep=0.1em]
    \item Calmer, large-cap assets receive higher capital allocations ($\approx 12\%\text{--}16\%$).
    \item Volatile small-cap assets receive lower capital allocations ($\approx 6\%\text{--}9\%$).
\end{itemize}
This equalizes the marginal risk contribution of each holding, dampening overall portfolio variance and protecting the portfolio during broad market corrections.

Target cash allocations $C_i = V_{\text{total}} \cdot w_i$ are converted into discrete share quantities:
\begin{equation}
    \text{Shares Target}_i = \left\lfloor \frac{C_i}{P_{i, t}} \right\rfloor
\end{equation}

% ==============================================================================
% 5. REBALANCING ENGINE & TRANSACTION FRICTION
% ==============================================================================
\section{5. Execution, Rebalancing Logic \& Transaction Friction}
\begin{itemize}[leftmargin=1.5em]
    \item \textbf{Review Frequency:} The algorithm reviews and rebalances the portfolio on the **first trading day of each calendar month** across 2021--2025 (60 rebalancing cycles).
    \item \textbf{Order of Execution:}
    \begin{enumerate}[label=\arabic*., itemsep=0.1em]
        \item \textit{Exits (SELL):} Any stock held in the previous month that is no longer in the new target portfolio is completely liquidated at market close. Cash is credited net of transaction costs.
        \item \textit{Target Adjustments (BUY / TRIM):} For retained and newly selected assets, the model computes $\Delta C_i = (V_{\text{total}} \cdot w_i) - \text{Current Value}_i$. If $\Delta C_i > 0$, shares are purchased; if $\Delta C_i < 0$, excess shares are trimmed.
    \end{enumerate}
    \item \textbf{Mandatory 0.1\% Transaction Fee Accounting:} On every single buy, sell, or trim execution, a flat 0.1\% ($0.0010$) friction penalty is deducted directly from cash:
    \begin{equation}
        \text{Fee} = 0.001 \times |\text{Traded Rupee Value}|
    \end{equation}
    \begin{equation}
        \text{Cash}_{\text{new}} = 
        \begin{cases}
            \text{Cash}_{\text{old}} + (\text{Proceeds} - \text{Fee}), & \text{for SELL / TRIM} \\
            \text{Cash}_{\text{old}} - (\text{Cost} + \text{Fee}), & \text{for BUY}
        \end{cases}
    \end{equation}
\end{itemize}

% ==============================================================================
% 6. EMPIRICAL PERFORMANCE & RESULTS
% ==============================================================================
\section{6. Backtest Performance Results (2021--2025)}
The backtest was executed over the full 5-year competition period (1 January 2021 to 31 December 2025) using historical data across the eligible universe. Table~\ref{tab:results} summarizes the complete quantitative performance.

\begin{table}[H]
\centering
\small
\caption{Comprehensive Portfolio Performance Diagnostics (2021--2025)}
\label{tab:results}
\vspace{0.4em}
\begin{tabular}{llr}
\toprule
\textbf{Category} & \textbf{Quantitative Metric} & \textbf{Strategy Value} \\
\midrule
\textbf{Capital Metrics} & Virtual Starting Capital & Rs.~1,00,00,000.00 \\
                         & Final Portfolio Value & Rs.~4,80,20,487.39 \\
                         & \textbf{Total Net PNL} & \textbf{+Rs.~3,80,20,487.39} \\
\midrule
\textbf{Return Statistics} & Absolute 5-Year Return & +380.20\% \\
                           & Annualized Compounded Growth (CAGR) & 36.86\% \\
\midrule
\textbf{Risk Diagnostics}  & Sharpe Ratio ($R_f = 0\%$) & 1.82 \\
                           & Maximum Drawdown (MDD) & -23.63\% \\
                           & Annualized Volatility & 18.42\% \\
\midrule
\textbf{Trading Statistics}& Trade Accuracy (Win Rate on Closed Rounds) & 66.82\% \\
                           & Gain-to-Loss Ratio & 1.14 \\
                           & Total Rebalance Executions & 758 orders \\
                           & Cumulative Transaction Costs Paid & Rs.~14,79,089.07 \\
\bottomrule
\end{tabular}
\end{table}

\begin{figure}[H]
    \centering
    \includegraphics[width=0.92\linewidth]{equity_curve.png}
    %\fbox{\rule{0pt}{2.0in} \rule{0.9\linewidth}{0pt}}
    \caption{Portfolio Capital Trajectory Compounding from Rs.~1.00 Cr to Rs.~4.80 Cr (2021--2025).}
    \label{fig:equity}
\end{figure}

% ==============================================================================
% 7. BENCHMARK COMPARISON & OUT-OF-SAMPLE DURABILITY
% ==============================================================================
\section{7. Benchmark Comparison \& Out-of-Sample Durability}
\subsection{7.1 Benchmark Outperformance}
Over the identical 2021--2025 horizon, broad market indices (Nifty 100 and Nifty 500) generated annualized returns (CAGR) between 14.5\% and 16.2\%. The strategy delivered an annualized return of \textbf{36.86\%}, generating over \textbf{20\% annualized alpha}. Crucially, this excess return was achieved with controlled downside risk: while broader mid/small-cap indices suffered peak drawdowns exceeding 30\% during mid-cycle corrections, the portfolio's maximum drawdown was contained at \textbf{-23.63\%}.

\subsection{7.2 Preparedness for Out-of-Sample Stress Testing (2026)}
The strategy avoids common quantitative failure modes:
\begin{itemize}[leftmargin=1.5em, itemsep=0.1em]
    \item \textbf{Zero Parameter P-Hacking:} The lookback windows (250-day momentum, 60-day correlation, 200-day SMA) and factor weights (60\% momentum / 40\% stability) are standard institutional parameters that were kept fixed across all 60 months.
    \item \textbf{Absence of Future Information Leakage:} Every rebalance decision on day $t$ utilizes strictly the information available on or before day $t$.
    \item \textbf{Regime Adaptability:} Because the selection rules and inverse-volatility engine are fully dynamic, when market leadership rotates across sectors or market-cap segments during the 2026 test, the model will systematically adapt and allocate capital to leading, de-correlated assets.
\end{itemize}


% ==============================================================================
% 6. PERFORMANCE & RISK EVALUATION METRIC FRAMEWORK
% ==============================================================================
\section{8. Performance \& Risk Evaluation Metric Framework}
To rigorously quantify portfolio profitability, risk-adjusted performance, and execution efficiency across the 2021--2025 backtest window, the strategy evaluates seven core quantitative metrics. Below are their formal mathematical definitions and financial interpretations.

\subsection{8.1 Capital Metrics: Total Net Profit and Loss (PNL) \& Absolute Return}
\begin{itemize}[leftmargin=1.5em]
    \item \textbf{Mathematical Definition:} Let $V_0 = \text{Rs.}~1,00,00,000$ denote the initial virtual corpus at $t = 0$, and $V_T$ denote the final portfolio liquidation value on date $T$ (31 December 2025). The Total Net PNL and Absolute Return ($R_{\text{total}}$) are defined as:
    \begin{equation}
        \text{Total Net PNL} = V_T - V_0
    \end{equation}
    \begin{equation}
        R_{\text{total}} = \frac{V_T - V_0}{V_0} = \frac{V_T}{V_0} - 1
    \end{equation}
    \item \textbf{Financial Logic \& Intuition:} Net PNL serves as the primary evaluation metric for the challenge, measuring the absolute monetary wealth generated after accounting for all market fluctuations, capital reallocation, and trade frictions. Absolute return standardizes this capital growth as an un-annualized percentage gain over the 5-year investment horizon.
\end{itemize}

\subsection{8.2 Compounded Growth: Compound Annual Growth Rate (CAGR)}
\begin{itemize}[leftmargin=1.5em]
    \item \textbf{Mathematical Definition:} Let $N = 5$ represent the total duration of the backtest in years (or $D$ elapsed calendar days). The Compound Annual Growth Rate is defined as the geometric average rate of return:
    \begin{equation}
        \text{CAGR} = \left( \frac{V_T}{V_0} \right)^{\frac{1}{N}} - 1 = \left( \frac{V_T}{V_0} \right)^{\frac{365.25}{D}} - 1
    \end{equation}
    \item \textbf{Financial Logic \& Intuition:} Unlike the arithmetic mean return (which severely distorts performance in the presence of volatility), CAGR captures multi-period geometric compounding. It reflects the constant annual rate of growth required for the capital base to compound from $V_0$ to $V_T$, heavily penalizing strategies with volatile path-dependent drawdowns.
\end{itemize}

\subsection{8.3 Downside Tail Risk: Maximum Drawdown (MDD)}
\begin{itemize}[leftmargin=1.5em]
    \item \textbf{Mathematical Definition:} Let $V_t$ be the portfolio valuation on daily timestamp $t \in [0, T]$, and let $M_t$ denote the running peak (high-water mark) up to day $t$:
    \begin{equation}
        M_t = \max_{\tau \in [0, t]} V_\tau
    \end{equation}
    The drawdown at time $t$, denoted $DD_t$, and the Maximum Drawdown ($\text{MDD}$) over the entire evaluation period are defined as:
    \begin{equation}
        DD_t = \frac{V_t - M_t}{M_t}, \qquad \text{MDD} = \min_{t \in [0, T]} DD_t
    \end{equation}
    \item \textbf{Financial Logic \& Intuition:} MDD quantifies the single worst peak-to-trough percentage loss an investor would have endured before a new equity peak was reached. In quantitative asset management, MDD is the premier metric for downside tail risk, stress-testing whether a strategy avoids catastrophic capital destruction during broad market corrections.
\end{itemize}

\subsection{8.4 Risk-Adjusted Efficiency: Sharpe Ratio ($R_f = 0\%$)}
\begin{itemize}[leftmargin=1.5em]
    \item \textbf{Mathematical Definition:} Let $r_{p, t} = \frac{V_t - V_{t-1}}{V_{t-1}}$ be the daily percentage portfolio returns across $K$ trading days, with mean $\bar{r}_p = \frac{1}{K}\sum_{t=1}^K r_{p, t}$ and sample standard deviation $\sigma_p = \sqrt{\frac{1}{K-1}\sum_{t=1}^K (r_{p, t} - \bar{r}_p)^2}$. Annualized return is $\mu_{\text{ann}} = 252 \cdot \bar{r}_p$ and annualized volatility is $\sigma_{\text{ann}} = \sigma_p \sqrt{252}$. Under the competition constraint of a zero risk-free rate ($R_f = 0\%$), the Sharpe Ratio is formulated as:
    \begin{equation}
        \text{Sharpe Ratio} = \frac{\bar{r}_p}{\sigma_p}\sqrt{252} = \frac{\mu_{\text{ann}}}{\sigma_{\text{ann}}}
    \end{equation}
    \item \textbf{Financial Logic \& Intuition:} The Sharpe ratio evaluates the efficiency of return generation per unit of total risk (variance). By fixing $R_f = 0\%$, the metric provides an unambiguous baseline to verify that high CAGR was achieved through steady, low-variance compounding rather than reckless, unhedged beta risk.
\end{itemize}

\subsection{8.5 Execution Diagnostics: Trade Accuracy \& Gain-to-Loss Ratio}
\begin{itemize}[leftmargin=1.5em]
    \item \textbf{Mathematical Definition:} Let $\mathcal{C}$ be the set of all liquidated round-trip positions over the 5-year simulation, where trade $k \in \mathcal{C}$ generated net profit $\text{PnL}_k = (\text{Proceeds}_k - \text{Fee}_k) - \text{Cost}_k$. The set of winning positions is $\mathcal{W} = \{k \in \mathcal{C} \mid \text{PnL}_k > 0\}$ and losing positions is $\mathcal{L} = \{k \in \mathcal{C} \mid \text{PnL}_k \le 0\}$.
    \begin{equation}
        \text{Trade Accuracy (Win Rate)} = \frac{|\mathcal{W}|}{|\mathcal{C}|} \times 100\%
    \end{equation}
    \begin{equation}
        \text{Gain-to-Loss Ratio} = \frac{\frac{1}{|\mathcal{W}|}\sum_{k \in \mathcal{W}} \text{PnL}_k}{\frac{1}{|\mathcal{L}|}\sum_{k \in \mathcal{L}} |\text{PnL}_k|}
    \end{equation}
    \item \textbf{Financial Logic \& Intuition:} 
    \begin{enumerate}[label=(\alph*), itemsep=0.1em]
        \item \textit{Trade Accuracy:} Measures the statistical hit rate of the dual-factor alpha signal in correctly identifying outperforming equities over a 1-month holding horizon.
        \item \textit{Gain-to-Loss Ratio:} Measures trade payoff asymmetry (profit factor). A strategy with an asymmetric Gain-to-Loss ratio ($>1.0$) remains profitable even through low-accuracy regimes by strictly cutting losers while allowing high-momentum winners to run.
    \end{enumerate}
\end{itemize}

\subsection{8.6 Friction \& Turnover Accounting}
\begin{itemize}[leftmargin=1.5em]
    \item \textbf{Mathematical Definition:} For each execution $m \in \{1, \dots, N_{\text{orders}}\}$ involving $|\Delta n_m|$ shares traded at market price $P_{m}$, the transaction fee is charged at flat 0.1\%:
    \begin{equation}
        \text{Friction Cost}_m = 0.001 \times (|\Delta n_m| \cdot P_{m})
    \end{equation}
    \begin{equation}
        \text{Cumulative Friction} = \sum_{m=1}^{N_{\text{orders}}} \text{Friction Cost}_m
    \end{equation}
    \item \textbf{Financial Logic \& Intuition:} Tracks real-world execution drag, bid-ask slippage, and brokerage overhead. Incorporating mandatory 0.1\% friction prevents high-turnover churn and ensures that reported net returns represent deployable institutional alpha.
\end{itemize}

\end{document}